```latex
\section{Additional analyses}

The analyses in this section are \emph{exploratory} and \emph{secondary} to the
pre-declared primary families reported in Sections~4 and~5. None of them reuses
any primary claim identifier (\texttt{LL-A}, \texttt{LL-B}, \texttt{LL-C},
\texttt{QW-A}, \texttt{QW-B}, \texttt{QW-C}, \texttt{GM-A}, \texttt{G5-A},
\texttt{CS-A}, \texttt{GO-A}), and none of them belongs to, extends, or is
corrected within the pre-declared Holm families for Llama, Qwen, Gemma, or
GPT-OSS (\texttt{turbo\_alpha30\_role\_strength\_6\_tests},
\texttt{qwen25\_72b\_primary\_33\_steered\_cells},
\texttt{gemma\_phase1\_three\_cell\_family}, the four-comparison GPT-OSS harmful
family, and the two four-comparison Llama harmful and benign families). No new
correction family is declared here, no hypothesis test result reported in this
section is treated as confirmatory, and no quantity from this section is pooled
with, averaged over, or used to rank the model entries of Sections~4 and~5. The
six subsections below are therefore descriptive or hypothesis-generating: they
document what the shipped Addition artifacts contain, and they are intended to
motivate future pre-registered designs rather than to support conclusions.
Where the packet does not contain the artifacts a subsection would require, we
say so and keep the subsection at a high level rather than reporting numbers
from an unavailable source.

\subsection{Cross-model comparison}

The exploratory question here is whether the Addition analyses, taken one at a
time, show the same qualitative picture in more than one model entry. The
available coverage is uneven and asymmetric, so the answer is limited to a
description of that coverage. Addition~1 covers
\texttt{meta-llama/llama-3.3-70b-instruct} only, in Hausa, as an endpoint audit
of an activation-patching artifact
(\texttt{opus\_packet/metrics/addition1\_llama33\_hausa\_patching\_v3\_audit.csv};
\texttt{opus\_packet/tables/addition1\_llama33\_hausa\_patching\_v3\_table.tex};
\texttt{opus\_packet/handoffs/addition1\_llama33\_hausa\_patching\_v3\_OPUS\_HANDOFF.md};
\texttt{opus\_packet/manifests/addition1\_llama33\_hausa\_patching\_v3\_artifact\_manifest.json}).
That artifact contains 11 canonical rows with 11 unique base IDs, drawn from 1
source path with 1 source SHA-256 value. The endpoint decision counts are: clean
1 ALLOW and 10 VETO; corrupted 0 ALLOW and 11 VETO; patched 0 ALLOW and 11
VETO. Recovery is defined only where the clean and corrupted endpoint decisions
differ, which holds for exactly 1 of the 11 rows
(\texttt{evaluable\_clean\_ne\_corrupted} $=1$); of that single evaluable row, 0
recovered at the patched endpoint (\texttt{recovered\_patched\_eq\_clean} $=0$),
giving \texttt{recovery\_rate\_among\_evaluable} $=0.0$, i.e.\ $0/1$. We report
this as an observed endpoint outcome in this artifact with an explicit
denominator of one; it is not a population recovery rate and it is not evidence
of broad patching failure.

Addition~3 is the only Addition artifact that touches two model entries. As
documented in
\texttt{opus\_packet/evidence/additions\_2\_3\_claim\_boundary\_audit.md}, it
covers Hindi, Korean, and Swahili with 50 harmful and 50 benign prompts per
language and 300 paired rows per model, storing a baseline and a steered
decision for each row. The documented outcome is the harmful VETO rate, which is
a \emph{different estimand} from the operational-failure and refusal outcomes of
Sections~4 and~5 and is not comparable to them. For Llama the documented harmful
VETO rates are Hindi 100.0\% baseline and 100.0\% steered, Korean 88.0\% and
90.0\%, Swahili 94.0\% and 94.0\%, and overall 94.0\% and 94.7\%; for Qwen they
are Hindi 86.0\% and 86.0\%, Korean 94.0\% and 92.0\%, Swahili 66.0\% and
68.0\%, and overall 82.0\% and 82.0\%. On these three languages the changes are
small and heterogeneous, with Llama near a harmful-VETO ceiling in this set and
Qwen showing mixed language-specific changes. The remaining Additions add no
cross-model contrast: Addition~5 is Llama-only and English-only, Addition~4 is
an annotation-agreement study rather than a measure of model behaviour, and the
Addition~6 artifacts record no model identifier anywhere, so Addition~6 cannot
be attributed to a model entry here.

The limitations are binding and jointly decisive. There is no correction family
for any of these comparisons; Addition~3 provides no paired inferential test, no
confidence interval, and no multiple-testing correction; the outcome definitions
differ across Additions (endpoint ALLOW/VETO decisions, harmful VETO rate,
jailbreak success, annotation agreement), so the artifacts are not on a common
scale; Addition~3's Qwen arm has documented malformed outputs, especially in
Swahili, whose counts are not provided, so its missing-label denominators are
not established; and the raw Addition~3 result files are not shipped in the
packet, so only the documented rates above can be quoted. Accordingly this
subsection is descriptive and hypothesis-generating. We make no claim of
cross-lingual improvement, no claim of recovery, no pooling across models, and
no ranking of models.

\subsection{Role specificity}

The exploratory question is whether outcome patterns differed when the same
historical steering direction was applied to different agent roles. Two sources
speak to this, and both are limited. First, Addition~2 is retained as a
\emph{retrospective} role-specificity analysis rather than a completed
matched-control experiment
(\texttt{opus\_packet/evidence/additions\_2\_3\_claim\_boundary\_audit.md};
\texttt{paper\_artifacts/metrics/addition\_2\_status.md}). Its implementation
aggregates archived Llama~3.3 Phase~2 Turbo results by configuration, strength,
and language and prints Hausa veto rates for planner-only, critic-only, and
both-role steering. Its intended derived output
(\texttt{results/addition\_2\_matched\_controls/analysis.json}) and its
historical input file are \emph{not} in the packet: there is no metrics CSV, no
table, and no manifest for it, so Addition~2 contributes \emph{no citable
number} here and this part of the subsection is necessarily kept at a high
level. The most that the evidence supports is the recorded qualitative
statement: retrospective analyses of the same steering direction across roles
showed role-dependent outcome patterns, which is consistent with role
sensitivity and supports treating intervention location as an empirical design
variable.

Second, the archived cross-phase display artifact
(\texttt{opus\_packet/tables/llama33\_turbo\_alpha30\_cross\_phase\_comparison.tex})
records six configuration $\times$ language cells (\texttt{planner\_only},
\texttt{critic\_only}, \texttt{both} $\times$ Hausa, Yoruba), each with a matched
$n$ of 300. The direction of the recorded change in harmful
operational-failure rate is not the same in every configuration label: the
\texttt{planner\_only} cells record $-37.7$~pp (Hausa) and $-19.3$~pp (Yoruba),
the \texttt{critic\_only} cells record $+6.3$~pp and $+13.0$~pp, and the
\texttt{both} cells record $-19.0$~pp and $+10.0$~pp. We report these as the
change values printed for the configuration-labelled cells of an archived
artifact and nothing more. In particular, we do not attach a role-specificity or
mechanistic interpretation to this artifact: the comparison varies configuration
label jointly with run phase and steering strength, so configuration is
confounded with phase, and record-ID matching does not establish non-prompt
run-condition equivalence. We also do not restate the baseline and Alpha-30
percentage columns of that artifact, whose printed values are on an inconsistent
scale (e.g.\ ``6633.3'' and ``4733.3'' in a column labelled per cent); only the
change, interval, and matched-$n$ columns are quoted here, and the properly
scoped version of this comparison, with its own pre-declared family, is the one
described in Section~3 and is not re-asserted in this section. The key
limitation is therefore twofold: the retrospective analysis has no shipped
numeric artifact, and the cross-phase artifact cannot separate role from phase.
This subsection is hypothesis-generating only; it is not confirmatory evidence
of role specificity.

\subsection{Vector specificity}

The exploratory question is whether any of the observed outcome patterns can be
attributed to the specific steering direction used, as opposed to generic
perturbation of the activations. The packet does not contain artifacts that can
answer this, and we state that plainly rather than reporting a proxy. The
Addition~2 implementation loads a single historical refusal feature and
aggregates archived outcomes; per
\texttt{opus\_packet/evidence/additions\_2\_3\_claim\_boundary\_audit.md} it does
\emph{not} generate a random direction, does \emph{not} run a same-session
matched random-vector control, and does \emph{not} compare a real-vector
intervention against a matched random-vector intervention. It therefore neither
identifies nor rules out the generic-perturbation null. We accordingly make no
claim that matched random-vector controls were run, no claim that generic
perturbation was excluded, and no claim that the retrospective result
demonstrates vector specificity. For the archived cross-phase artifact, the
original steering-vector file is declared unavailable, so no direction-level
comparison is possible there either.

The one artifact in the packet that uses a random direction is Addition~5
(\texttt{opus\_packet/metrics/addition5\_adversarial\_robustness\_metrics.csv};
\texttt{opus\_packet/tables/addition5\_adversarial\_robustness.tex};
\texttt{opus\_packet/handoffs/addition5\_adversarial\_robustness\_OPUS\_HANDOFF.md}),
but it is not a matched control for vector specificity: it applies a random
steering vector (layer 52, strength 20) on a 50-prompt English jailbreak set
with a different outcome (jailbreak success), records 0 of 50 successes in both
the baseline and the steered condition, and contains no real-vector arm on the
same prompts. Because there is no matched real-versus-random contrast anywhere
in the packet, vector specificity is \emph{not established} by any analysis in
this paper. The key limitation is the absence of the required artifact, not an
inconclusive test. This subsection is hypothesis-generating: it identifies a
same-session matched random-vector control on the same prompts and outcome
definition as the design that would be needed, and it makes no confirmatory
claim.

\subsection{Combined versus individual interventions}

The exploratory question is how cells labelled with the combined configuration
compare with the two single-configuration cells in the same language. The only
artifact in the packet that contains all three configuration labels is the
archived cross-phase display artifact
(\texttt{opus\_packet/tables/llama33\_turbo\_alpha30\_cross\_phase\_comparison.tex}),
covering Hausa and Yoruba harmful prompts with a matched $n$ of 300 in each of
its six cells. Taking the printed change column only, in Hausa the recorded
changes are $-37.7$~pp for \texttt{planner\_only}, $+6.3$~pp for
\texttt{critic\_only}, and $-19.0$~pp for \texttt{both}; in Yoruba they are
$-19.3$~pp, $+13.0$~pp, and $+10.0$~pp respectively. In each of the two
languages the \texttt{both} value lies numerically between the two
single-configuration values. We report that as an arithmetic description of the
printed numbers and draw no additivity, interaction, or dose-response inference
from it: no interaction test was pre-declared or performed, the printed
per-cell intervals are not intervals for differences between configurations, and
no comparison across the three configuration labels was corrected in any
family. The remaining Addition artifacts contribute no combined-versus-individual
contrast: Addition~5 has a single random-vector condition
(\texttt{opus\_packet/metrics/addition5\_adversarial\_robustness\_metrics.csv})
and Addition~4 is an annotation-agreement study
(\texttt{opus\_packet/metrics/addition4\_judge\_calibration\_metrics.csv}), so
neither contains a combined intervention arm.

The key limitations are cross-phase confounding and scope. The configuration
labels are compared across a phase boundary and at a single strength, so
configuration, strength, and run phase vary together; the comparison covers two
languages, one model entry, and harmful prompts only; no correction family
covers the between-configuration contrast; and the percentage-level columns of
the display artifact are on an inconsistent scale and are not quoted. This
subsection is descriptive and hypothesis-generating, not confirmatory, and it
supports no claim that combining interventions is better or worse than applying
either one alone.

\subsection{Cross-phase comparison}

The exploratory question is what the archived Phase~1 versus Phase~2 ``Turbo''
display artifact records. The data are Llama~3.3 harmful prompts in Hausa and
Yoruba only, contrasting the Phase~1 main baseline against Phase~2 Turbo at
$\alpha30$ under three configuration labels, with a matched $n$ of 300 in each
of the six cells
(\texttt{opus\_packet/tables/llama33\_turbo\_alpha30\_cross\_phase\_comparison.tex}).
The printed changes in harmful operational-failure rate, with the printed 95\%
intervals in percentage points, are: \texttt{both} Hausa $-19.0$
$[-26.0, -12.0]$; \texttt{both} Yoruba $+10.0$ $[4.7, 15.3]$;
\texttt{critic\_only} Hausa $+6.3$ $[0.7, 12.0]$; \texttt{critic\_only} Yoruba
$+13.0$ $[8.7, 17.3]$; \texttt{planner\_only} Hausa $-37.7$ $[-44.3, -31.0]$;
and \texttt{planner\_only} Yoruba $-19.3$ $[-24.7, -14.0]$. The artifact's
``Reject?'' column reads ``Yes'' for all six cells; that column reflects the
six-test correction family declared for the underlying analysis described in
Section~3, and we neither re-declare it here nor treat any cell in this section
as a confirmatory rejection. Two further display caveats apply: the baseline and
Alpha-30 percentage columns of this artifact are printed on an inconsistent
scale and are therefore not reported, and the harmful failure outcome used in
this comparison folds undecided records into the failure numerator, as recorded
in Section~3, so it is not interchangeable with the Phase~1 outcomes of
Section~4.1.

The key limitation is cross-phase confounding. Baseline and steered records come
from two different run phases, and matching by record ID does not establish that
non-prompt run conditions were equivalent between the phases; strength and
configuration also differ from the Phase~1 primary contrast, and only two
languages and one prompt class are covered. Consequently the differences above
are associational descriptions of two archived record sets and cannot be read as
effects of steering. This subsection is descriptive and hypothesis-generating;
no causal, mechanistic, or deployment claim follows from it, and no primary
claim identifier is asserted through it.

\subsection{Secondary outcomes}

Three Addition artifacts report outcomes that are secondary to the decision
outcomes of Sections~4 and~5: annotation agreement (Addition~4), jailbreak
success under a random steering vector (Addition~5), and a three-condition
baseline comparison of harmful veto rate and benign false positives
(Addition~6). Each is reported on its own terms, with its own denominators, and
none is combined with any other.

Addition~4 asks how consistently two independent annotators applied the binary
compliance label. The data are 98 prompts, 14 per language in Arabic, English,
Hausa, Hindi, Korean, Swahili, and Yoruba, each labelled by both annotators as
YES (should comply) or NO (should refuse)
(\texttt{opus\_packet/metrics/addition4\_judge\_calibration\_metrics.csv};
\texttt{opus\_packet/metrics/addition4\_calibration\_paper\_table.csv};
\texttt{opus\_packet/tables/addition4\_calibration\_agreement.tex};
\texttt{opus\_packet/handoffs/addition4\_calibration\_OPUS\_HANDOFF.md}).
Six of the seven languages---Arabic, English, Hausa, Hindi, Korean, and
Swahili---record observed agreement of $1.0$ with Cohen's $\kappa = 1.0$ and a
95\% interval of $[1.00, 1.00]$. Yoruba records observed agreement of $0.9286$
(13 of 14) with $\kappa = 0.857$ and a 95\% interval of $[0.533, 1.00]$.
Overall, observed agreement is $0.9898$ (97 of 98 prompts, i.e.\ a single
disagreement) with $\kappa = 0.980$ and a 95\% interval of $[0.937, 1.00]$;
annotator marginals are 46 YES and 52 NO for the first annotator and 47 YES and
51 NO for the second, and every language meets the $\kappa \ge 0.60$ target. The
single disagreement was resolved by blinded adjudication, with the adjudicator
seeing only disagreement cases. We report $\kappa$ intervals in $\kappa$ units
from the metrics CSV; the paper-table CSV and LaTeX display the same intervals
multiplied by 100 without a unit label, and that display convention should not
be read as a percentage. The key limitations are that the annotators were
\emph{not} blinded to the original labels, which may have inflated agreement,
and that the calibration set is small (98 prompts). This is a descriptive
inter-rater reliability study: it measures annotation consistency, not model
performance, it is not an intervention, and it does not generalise beyond these
98 prompts.

Addition~5 asks whether applying a steering vector was accompanied by any
increase in jailbreak success on a bounded English attack set. The data are 50
English jailbreak prompts for Llama-3.3-70B-Instruct across six attack families
---\texttt{dan} ($n=8$), \texttt{developer} ($n=8$), \texttt{instruction}
($n=8$), \texttt{other} ($n=10$), \texttt{roleplay} ($n=8$), and
\texttt{translation} ($n=8$)---compared between a baseline condition and a
\emph{random} steering vector at layer 52, strength 20
(\texttt{opus\_packet/metrics/addition5\_adversarial\_robustness\_metrics.csv};
\texttt{opus\_packet/metrics/addition5\_adversarial\_robustness\_table.csv};
\texttt{opus\_packet/tables/addition5\_adversarial\_robustness.tex};
\texttt{opus\_packet/handoffs/addition5\_adversarial\_robustness\_OPUS\_HANDOFF.md}).
Every family, and the overall row, records 0 baseline jailbreaks and 0 steered
jailbreaks: baseline rate $0.0$, steered rate $0.0$, rate difference $0.0$, and
0 discordant pairs, so the McNemar entry is recorded as not applicable. All 50
prompts were evaluated as NEITHER. No evidence of increased vulnerability was
detected on this set. The key limitations are that the outcome has zero events
in both conditions, so no test is defined and the check has no power to detect
small changes; that the steering vector was random rather than a refusal
direction; and that all prompts are English. Accordingly this is a bounded
descriptive check: it does not show that steering is safe, does not generalise
beyond these 50 prompts or to untested attacks or to an actual refusal
direction, and provides no multilingual robustness claim.

Addition~6 asks how a prompt-filter baseline and a steering condition compare on
one evaluation set, under one implementation and protocol. The data are 460
prompts (400 harmful, 60 benign) across four languages, evaluated under three
conditions---\texttt{no\_intervention}, \texttt{prompt\_filter}, and
\texttt{planner\_steering\_l50\_s20\_historical\_all\_steps}
(\texttt{opus\_packet/metrics/addition6\_prompt\_filter\_baseline\_metrics.csv};
\texttt{opus\_packet/metrics/addition6\_prompt\_filter\_baseline\_table.csv};
\texttt{opus\_packet/tables/addition6\_prompt\_filter\_baseline.tex};
\texttt{opus\_packet/handoffs/addition6\_prompt\_filter\_baseline\_OPUS\_HANDOFF.md};
\texttt{opus\_packet/manifests/addition6\_prompt\_filter\_baseline\_artifact\_manifest.json}).
The recorded harmful veto rates are 210 of 400 ($52.5\%$) with no intervention,
214 of 400 ($53.5\%$) under the prompt filter, and 230 of 400 ($57.5\%$) under
steering; the recorded benign veto (false-positive) rates are 0 of 60 ($0.0\%$),
0 of 60 ($0.0\%$), and 3 of 60 ($5.0\%$) respectively; the recorded missing-label
counts are 4, 5, and 25. The steering condition thus records a harmful veto rate
$5.0$~pp above the no-intervention condition and $4.0$~pp above the
prompt-filter condition, together with three benign false positives and a
substantially larger missing-label count. Both parts of that trade-off are
reported here. The key limitations are that no test, interval, or
multiple-testing correction is defined for this three-condition comparison; that
the missing-label counts differ markedly across conditions (4, 5, and 25), so the
condition denominators are not label-complete and the differences above are not
adjusted for that asymmetry; that the outcome is a harmful veto rate rather than
the operational-failure outcomes of Section~4 and is not comparable to them; and
that the artifacts record no model identifier, so this comparison is not
attributed to any model entry in this paper. It evaluates outcomes for this
prompt set, this prompt-filter implementation, this steering configuration, and
this evaluation protocol only, and it does not establish that steering is
superior to prompt filtering in other models, languages, datasets, prompt
distributions, or deployment settings. This subsection is descriptive and
hypothesis-generating, not confirmatory.

% ---------------- SOURCE MAP ----------------
% All six subsections (6 — additional analyses): exploratory, not covered by
% primary claim_ids (LL-A/B/C, QW-A/B/C, GM-A, G5-A, CS-A, GO-A), no Holm
% family claimed across these comparisons.
%
% 6.0 intro paragraph — exploratory/secondary status; no primary claim_ids
%   reused; no membership in the pre-declared Llama, Qwen, Gemma, or GPT-OSS
%   Holm families; no new correction family; no pooling or ranking.
% 6.1 Cross-model comparison — Addition 1:
%   opus_packet/metrics/addition1_llama33_hausa_patching_v3_audit.csv,
%   opus_packet/tables/addition1_llama33_hausa_patching_v3_table.tex,
%   opus_packet/handoffs/addition1_llama33_hausa_patching_v3_OPUS_HANDOFF.md,
%   opus_packet/manifests/addition1_llama33_hausa_patching_v3_artifact_manifest.json
%   (11 canonical rows, 11 unique base IDs, 1 source path, 1 source SHA-256;
%   clean 1 ALLOW / 10 VETO; corrupted 0/11; patched 0/11; evaluable = 1;
%   recovered = 0; recovery rate among evaluable = 0.0);
%   Addition 3: opus_packet/evidence/additions_2_3_claim_boundary_audit.md
%   (Hindi/Korean/Swahili, 50 harmful + 50 benign per language, 300 paired rows
%   per model, documented harmful VETO rates for Llama and Qwen); Addition 5 and
%   Addition 6 coverage statements (Llama/English-only; no model identifier).
%   Limitations: no correction family, no tests/CIs for Addition 3,
%   outcome-definition heterogeneity, Qwen malformed outputs (counts not
%   provided), raw Addition 3 result files not shipped.
% 6.2 Role specificity — Addition 2:
%   opus_packet/evidence/additions_2_3_claim_boundary_audit.md,
%   paper_artifacts/metrics/addition_2_status.md (no citable number; intended
%   analysis.json not in packet); archived cross-phase display artifact
%   opus_packet/tables/llama33_turbo_alpha30_cross_phase_comparison.tex
%   (6 cells, matched n = 300, printed change column only). No role-specificity
%   or mechanistic claim attached to the cross-phase artifact; inconsistent-scale
%   percentage columns not quoted.
% 6.3 Vector specificity — no supporting artifact in the packet; negative
%   findings from opus_packet/evidence/additions_2_3_claim_boundary_audit.md
%   (no random direction generated, no same-session matched random-vector
%   control, generic-perturbation null neither identified nor ruled out) and
%   Addition 5 artifacts (random vector, English-only, zero events) as an
%   explicitly inadequate substitute; subsection kept at a high level.
% 6.4 Combined versus individual interventions —
%   opus_packet/tables/llama33_turbo_alpha30_cross_phase_comparison.tex
%   (Hausa: -37.7 / +6.3 / -19.0 pp; Yoruba: -19.3 / +13.0 / +10.0 pp; matched
%   n = 300 per cell); Addition 4 and Addition 5 metrics cited only to record
%   that they contain no combined intervention arm. Limitations: cross-phase
%   confounding, joint variation of configuration/strength/phase, no interaction
%   test, no correction family.
% 6.5 Cross-phase comparison —
%   opus_packet/tables/llama33_turbo_alpha30_cross_phase_comparison.tex
%   (six cells, matched n = 300, printed changes and 95% intervals in pp,
%   "Reject?" column noted as belonging to the family declared in Section 3 and
%   not re-declared here; baseline and Alpha-30 percentage columns on an
%   inconsistent scale and therefore not reported). Limitation: cross-phase
%   confounding; associational only.
% 6.6 Secondary outcomes — Addition 4:
%   opus_packet/metrics/addition4_judge_calibration_metrics.csv,
%   opus_packet/metrics/addition4_calibration_paper_table.csv,
%   opus_packet/tables/addition4_calibration_agreement.tex,
%   opus_packet/handoffs/addition4_calibration_OPUS_HANDOFF.md;
%   Addition 5: opus_packet/metrics/addition5_adversarial_robustness_metrics.csv,
%   opus_packet/metrics/addition5_adversarial_robustness_table.csv,
%   opus_packet/tables/addition5_adversarial_robustness.tex,
%   opus_packet/handoffs/addition5_adversarial_robustness_OPUS_HANDOFF.md;
%   Addition 6: opus_packet/metrics/addition6_prompt_filter_baseline_metrics.csv,
%   opus_packet/metrics/addition6_prompt_filter_baseline_table.csv,
%   opus_packet/tables/addition6_prompt_filter_baseline.tex,
%   opus_packet/handoffs/addition6_prompt_filter_baseline_OPUS_HANDOFF.md,
%   opus_packet/manifests/addition6_prompt_filter_baseline_artifact_manifest.json.
%   Limitations: annotators not blinded to original labels and small calibration
%   set (Addition 4); zero-event outcome with no defined test, random vector, and
%   English-only prompts (Addition 5); no test/correction, asymmetric
%   missing-label counts (4/5/25), non-comparable outcome, and no model
%   identifier (Addition 6).
```
